%-------------------------------------------------------------------- Comente: venta atada vs empaquetamiento
\needspace{8\baselineskip}
\tema{Comente: venta atada vs empaquetamiento}
\question \textbf{Comente.} Una firma que es monopolio en dos productos, con el objetivo de maximizar sus beneficios, siempre preferirá realizar una venta atada por sobre una venta empaquetada (empaquetamiento) de sus productos, ya que, al limitar al máximo las opciones de las y los consumidores, es capaz de extraerles el máximo excedente posible.
\needspace{5\baselineskip}
\begin{solution}
Falso. Preferirá el empaquetamiento por sobre la venta atada, ya que en el primero tiene más instrumentos para poder llevarse un mayor excedente de los consumidores.

\end{solution}

%-------------------------------------------------------------------- Comente: beneficio consumidor
\needspace{8\baselineskip}
\tema{Comente: ¿benefician al consumidor la venta atada y el empaquetamiento?}
\question \textbf{Comente.} La venta atada y el empaquetamiento son prácticas que benefician al consumidor.
\needspace{5\baselineskip}
\begin{solution}
Incierto. Efectivamente hay un conjunto de consumidores que se ven beneficiados: quienes habrían comprado ambos bienes incluso sin una venta atada o sin empaquetamiento. Sin embargo, existe otro grupo de consumidores que se ven perjudicados por estas prácticas, en específico quienes en el caso inicial solo habrían comprado uno de los dos bienes. En el caso de venta atada, parte de dicho grupo de consumidores se verá obligado a comprar el paquete de ambos bienes y pagar un precio mayor que el precio inicial. En el caso de empaquetamiento, aquel grupo de consumidores tiene la opción de solo comprar el bien que más valoren, pero a un precio mayor que el precio inicial.

\end{solution}

%-------------------------------------------------------------------- Comente: packs
\needspace{8\baselineskip}
\tema{Comente: ofertas en formato «pack»}
\question \textbf{Comente.} Las ofertas en formato «pack» que realizan las firmas no tienen racionalidad económica, pues tendrían mayor beneficio vendiendo los productos individualmente.
\needspace{5\baselineskip}
\begin{solution}
Falso. Existen razones económicas para realizar estas ofertas. El empaquetamiento otorga una mayor flexibilidad a la empresa, pues dispone de más instrumentos u ofertas con que segmentar a los consumidores. Otras razones son la economía de escala producto de la eficiencia del bien (por ejemplo, telecomunicaciones).

\end{solution}

%-------------------------------------------------------------------- Comente: per se
\needspace{8\baselineskip}
\tema{Comente: la venta atada como práctica anticompetitiva \emph{per se}}
\question \textbf{Comente.} La venta conjunta de productos, en particular la venta atada, al imponer que la adquisición de un producto deba ser en conjunto con otro(s) producto(s), sin opción de comprarlo por separado, siempre es perseguida y sancionada por las agencias de libre competencia, dado que es una práctica anticompetitiva \emph{per se}.
\needspace{5\baselineskip}
\begin{solution}
Las ventas conjuntas de productos pueden tener distintas justificaciones. En primer lugar, pueden ocurrir por un tema de eficiencia de costos para la firma. En segundo lugar, puede ser una manifestación del poder de mercado de la empresa en un mercado oligopólico. En tercer lugar, puede ser el resultado de una práctica anticompetitiva por parte de la firma para excluir competidores. Debido a las distintas explicaciones que puede tener una venta conjunta, dicha estrategia no es una práctica anticompetitiva \emph{per se}, sino que está sujeta a la \emph{regla de la razón}. Es decir, se debe evaluar caso a caso con el objetivo de dilucidar si dicha práctica tiene una justificación como una estrategia competitiva o bien como una práctica anticompetitiva. Por lo tanto, la afirmación es falsa.

\end{solution}

%-------------------------------------------------------------------- Comente: crítica de Chicago
\needspace{8\baselineskip}
\tema{Comente: crítica de Chicago (bienes complementarios)}
\question \textbf{Comente.} Existen 2 bienes perfectamente complementarios: $A$ y $B$. Una firma produce ambos bienes, donde en el mercado del bien $A$ es un monopolio y el del bien $B$ es de competencia perfecta. Además, todos los consumidores tienen la misma valoración por los bienes y todas las firmas tienen los mismos costos para producir el bien $B$. En esta situación, el monopolio tiene incentivos a desplegar venta atada para monopolizar el otro mercado y aumentar sus beneficios.
\needspace{5\baselineskip}
\begin{solution}
Falso. Si operara como monopolio en el mercado del bien $B$, no obtiene ganancia alguna. Esto se conoce como la \emph{crítica de Chicago}, la cual señala que es irrelevante monopolizar 2 bienes en vez de 1 cuando estos son perfectos complementos. Esto es posible porque los individuos solo obtienen utilidad al consumir simultáneamente ambos bienes (el monopolista ya extrae todo el excedente con el precio del bien $A$).

\end{solution}

%-------------------------------------------------------------------- Comente: fallo créditos
\needspace{8\baselineskip}
\tema{Comente: fallo sobre ventas atadas y empaquetamiento en créditos}
\question \textbf{Comente.} Hace unos años, un fallo penalizó las ventas atadas en el mercado de los créditos. Ese mismo fallo determinó que no había problemas con el empaquetamiento, ya que esas prácticas eran beneficiosas para los consumidores.
\needspace{5\baselineskip}
\begin{solution}
El empaquetamiento es una práctica donde se cobra un precio descontado si se consumen dos productos. A pesar de ello, esta estrategia permite cobrarle más a los consumidores de mayor valoración, por lo que es una forma de discriminar precios. Es así como no es claro si efectivamente beneficia a los consumidores: los consumidores de mayor valoración por un bien en particular y baja por el otro se verán perjudicados, ya que tendrán un precio mayor al caso sin empaquetamiento.

\end{solution}

%-------------------------------------------------------------------- Comente: venta atada y entrada
\needspace{8\baselineskip}
\tema{Comente: venta atada como disuasión de entrada}
\question \textbf{Comente.} Ventas atadas es un mecanismo de disuasión de entrada. Verdadero o falso. Justifique.
\needspace{5\baselineskip}
\begin{solution}
Verdadero. Además de extraer excedente al consumidor, usar ventas atadas es una forma de bloquear la entrada de competidores. El entrante obtiene menores beneficios si el incumbente utiliza ventas atadas, lo que reduce las probabilidades de entrada.

\end{solution}

%-------------------------------------------------------------------- Venta atada y empaquetamiento (uniforme)
\needspace{8\baselineskip}
\tema{Venta separada, venta atada y empaquetamiento (valoraciones uniformes)}
\question Un monopolista vende dos productos independientes $X$ e $Y$. El costo marginal de producción de cada bien es cero. Existe un conjunto de consumidores que normalizamos a una cantidad total de 1, los que valoran los bienes en $V_x$ y $V_y$, respectivamente. Las valoraciones se distribuyen uniformemente en el espacio $[0,1]\times[0,1]$.
\begin{parts}
\part Obtenga la función de beneficios a maximizar y grafique las demandas si es que el monopolista decide vender los productos por separado. Adicionalmente, señale en el espacio $V_x$--$V_y$ qué individuos consumen de cada tipo de bien por separado, de ambos bienes y de ningún bien.
\needspace{5\baselineskip}
\begin{solution}
El monopolista maximiza
\[
  \max_{p_x,p_y}\ \pi = (1 - p_x)p_x + (1 - p_y)p_y ,
\]
por lo que $p_x = p_y = \frac{1}{2}$ y $\pi = \frac{1}{2}$. Gráficamente, compran el bien $X$ quienes tienen $V_x \geq p_x$, y el bien $Y$ quienes tienen $V_y \geq p_y$:
\begin{center}
\cuadrado{
  \draw[dashed] (0.5,0) -- (0.5,1);
  \draw[dashed] (0,0.5) -- (1,0.5);
  \node at (0.25,0.25) {Ninguno};
  \node at (0.75,0.25) {Solo $X$};
  \node at (0.25,0.75) {Solo $Y$};
  \node at (0.75,0.75) {$X$ e $Y$};
  \node[below] at (0.5,0) {$p_x$};
  \node[left] at (0,0.5) {$p_y$};
}
\end{center}

\end{solution}

\part Obtenga la función de beneficios a maximizar y grafique las demandas si es que el monopolista decide vender los productos mediante venta atada. Adicionalmente, señale en el espacio $V_x$--$V_y$ qué individuos consumen de cada tipo de bien por separado, de ambos bienes y de ningún bien. En el gráfico muestre qué consumidores se ven beneficiados y qué consumidores se ven perjudicados con esta estrategia con respecto al anterior. ¿Qué nuevos individuos consumen con respecto al caso anterior? ¿Qué individuos dejan de consumir con respecto al caso anterior?
\needspace{5\baselineskip}
\begin{solution}
Si el monopolista decide vender los productos mediante venta atada, solo se vende el paquete $X+Y$ a un precio $p_{xy}$ y compran quienes cumplen $V_x + V_y \geq p_{xy}$. La función a maximizar es
\[
\begin{gathered}
  \max_{p_{xy}}\ \pi = \Bigl(1 - \frac{p_{xy}^2}{2}\Bigr)p_{xy} \\
  \;\Rightarrow\; 1 - \frac{3}{2}p_{xy}^2 = 0 \\
  \;\Rightarrow\; p_{xy} = \sqrt{\tfrac{2}{3}} \approx 0{,}82, \\
  \qquad \pi \approx 0{,}544 > 0{,}5 .
\end{gathered}
\]
\begin{center}
\cuadrado{
  \draw[dashed] (0.5,0) -- (0.5,1);
  \draw[dashed] (0,0.5) -- (1,0.5);
  \draw[thick] (0,0.8165) -- (0.8165,0);
  \node[font=\scriptsize,align=center] at (0.3,0.3) {Ninguno};
  \node[font=\scriptsize,align=center] at (0.72,0.72) {$X$ e $Y$\\(paquete)};
}
\end{center}
Con respecto al caso anterior (líneas punteadas):
\begin{itemize}
  \item \emph{Beneficiados:} quienes compraban ambos bienes por separado ($V_x, V_y \geq \frac12$) y que ahora pagan $0{,}82$ en vez de $1$.
  \item \emph{Nuevos consumidores:} quienes tienen $V_x, V_y < \frac12$ pero $V_x + V_y \geq 0{,}82$ (compran el paquete y antes no compraban nada).
  \item \emph{Dejan de consumir:} quienes antes compraban solo un bien (alta valoración por uno, baja por el otro) y cumplen $V_x + V_y < 0{,}82$. Los que sí siguen comprando pagan ahora más que antes por el bien que valoran y además se ven obligados a comprar el otro: están perjudicados.
\end{itemize}

\end{solution}

\part Obtenga la función de beneficios a maximizar y grafique las demandas si es que el monopolista decide vender los productos mediante empaquetamiento. Adicionalmente, señale en el espacio $V_x$--$V_y$ qué individuos consumen de cada tipo de bien por separado, de ambos bienes y de ningún bien. ¿Cuál de los tres escenarios le generará mayores beneficios al monopolista? Fundamente su respuesta.
\needspace{5\baselineskip}
\begin{solution}
En el caso en el que el monopolista aplique empaquetamiento (precios $p_x$, $p_y$ y $p_{xy}$), debe considerar las restricciones de participación y autoselección para cada uno de los precios.

\emph{Demanda que solo consume $X$.} Participación: $V_x - p_x \geq 0$. Autoselección: $V_x - p_x \geq V_x + V_y - p_{xy}$, es decir, $V_y \leq p_{xy} - p_x$. Juntando ambas:
\[
  D_x = (1 - p_x)(p_{xy} - p_x), \qquad \text{y de igual forma} \qquad D_y = (1 - p_y)(p_{xy} - p_y).
\]
\emph{Demanda por el paquete.} Participación: $V_x + V_y \geq p_{xy}$. Autoselección: $V_x \geq p_{xy} - p_y$ y $V_y \geq p_{xy} - p_x$. Juntando las restricciones:
\[
  D_{xy} = (1 - p_{xy} + p_y)(1 - p_{xy} + p_x) - \tfrac{1}{2}(p_x + p_y - p_{xy})^2 .
\]
Por lo tanto, la función a maximizar es
\[
\begin{gathered}
  \max\ \pi = p_x(1 - p_x)(p_{xy} - p_x) + p_y(1 - p_y)(p_{xy} - p_y) \\ + p_{xy}\Bigl[(1 - p_{xy} + p_y)(1 - p_{xy} + p_x) - \tfrac{1}{2}(p_x + p_y - p_{xy})^2\Bigr].
\end{gathered}
\]
Resolviendo numéricamente (caso simétrico) se obtiene $p_x = p_y \approx \frac{2}{3}$ y $p_{xy} \approx 0{,}86$, con $\pi \approx 0{,}549$. Gráficamente:
\begin{center}
\cuadrado{
  \draw (0.6667,0) -- (0.6667,0.1967) -- (1,0.1967);
  \draw (0,0.6667) -- (0.1967,0.6667) -- (0.1967,1);
  \draw (0.1967,0.6667) -- (0.6667,0.1967);
  \node[font=\scriptsize] at (0.28,0.28) {Ninguno};
  \node[font=\scriptsize] at (0.83,0.09) {Solo $X$};
  \node[font=\scriptsize] at (0.1,0.9) {Solo $Y$};
  \node[font=\scriptsize,align=center] at (0.7,0.7) {Paquete\\$X$ e $Y$};
}
\end{center}
El escenario que mayores beneficios reporta al monopolista es este último ($0{,}549 > 0{,}544 > 0{,}5$), debido a que, al tener más instrumentos disponibles, el monopolista puede crear mejores ofertas para sus distintos grupos de consumidores.

\end{solution}
\end{parts}

%-------------------------------------------------------------------- Bundling con tabla
\needspace{8\baselineskip}
\tema{Bundling puro y mixto con tres tipos de consumidores}
\question Considere la siguiente tabla, que contiene las valoraciones para los bienes $X$ e $Y$. Suponga que el costo marginal de $X$ es 1 y el costo marginal de $Y$ también es 1.
\begin{center}
\begin{tabular}{lcc}
  & Producto $X$ & Producto $Y$ \\ \hline
  Consumidor tipo 1 & 4 & 3 \\
  Consumidor tipo 2 & 3 & 3 \\
  Consumidor tipo 3 & 0 & 4
\end{tabular}
\end{center}
\begin{parts}
\part Encuentre el precio de monopolio óptimo en caso de ventas atadas o \emph{bundling} puro, es decir, precio de $X+Y$.
\needspace{5\baselineskip}
\begin{solution}
Bajo ventas atadas el monopolio vende $X+Y$ o no vende nada.
\begin{itemize}
  \item Si $P(X+Y) = 7$, solo compra el consumidor tipo 1 y los beneficios son $\pi = 7 - 1 - 1 = 5$.
  \item Si $P(X+Y) = 6$, compran los consumidores tipo 1 y 2, y los beneficios son $\pi = 2\cdot(6 - 1 - 1) = 8$.
  \item Si $P(X+Y) = 4$, compran los 3 tipos de consumidores, y los beneficios son $\pi = 3\cdot(4 - 1 - 1) = 6$.
\end{itemize}
Luego, la mejor estrategia es $P(X+Y) = 6$ con beneficios $\pi = 8$.

\end{solution}

\part Encuentre el precio de monopolio óptimo en caso de empaquetamiento o \emph{bundling} mixto, es decir, vender por separado y en conjunto. ¿Cuál de las estrategias, ventas atadas puras o mixtas, conlleva el mayor beneficio para el monopolista?
\needspace{5\baselineskip}
\begin{solution}
Se maximizan los beneficios vendiendo el paquete a un precio $P(X+Y) = 6$ y el bien $Y$ a $P(Y) = 4$. Note que el precio de $X$ solo es irrelevante, pues ningún tipo de consumidor compra el bien $X$ solo (se puede fijar, por ejemplo, un precio alto). El consumidor tipo 3 compra solo $Y$ (que valora en 4), y los tipos 1 y 2 compran el paquete. Los beneficios son
\[
  \pi = 2\cdot(6 - 1 - 1) + (4 - 1) = 11 > 8 .
\]
Por lo tanto, el \emph{bundling} mixto conlleva el mayor beneficio para el monopolista.

\end{solution}
\end{parts}

%-------------------------------------------------------------------- Entrada y venta atada
\needspace{8\baselineskip}
\tema{Entrada y venta atada: Agro S.A. y Siembra S.A.}
\question La empresa agrícola Agro S.A. siembra, cosecha y vende manzanas y naranjas, y actualmente posee el monopolio de ambos mercados. Se sabe además que las demandas que se enfrentan en estos mercados son de la forma $Q(P) = 4 - P$ en cada mercado (se puede ver así que las máximas disposiciones a pagar son de 4 en cada mercado). El costo marginal de producción tanto para manzanas como para naranjas es de 1. Agro se ha enterado de que Siembra S.A. está evaluando entrar al mercado de las manzanas, quien de hacerlo lo haría con un costo marginal igual a cero, pero debe enfrentar un costo de entrada $F = 2$ para hacerlo.
\begin{parts}
\part Evalúe si una estrategia en donde Agro S.A. vende una canasta de manzanas y naranjas podría disuadir la entrada de Siembra S.A.
\needspace{5\baselineskip}
\begin{solution}
El \emph{timing} del juego es:
\begin{itemize}
  \item $T=1$: Siembra decide si realiza su entrada al mercado o no.
  \item $T=2$: Agro y Siembra compiten en precios.
\end{itemize}
Resolviendo por inducción hacia atrás, en $T=2$ hay 2 posibles casos.

\textbf{Caso 1: Agro acomoda la entrada (no hay venta atada).} Siembra entra en el mercado de las manzanas quedándose con éste, al ser más eficiente y al estar compitiendo en precios. Agro sigue como monopolio en el mercado de las naranjas. Calculamos los beneficios de cada empresa:
\[
  \max_Q\ \pi_A = Q(4 - Q) - Q \;\Rightarrow\; 4 - 2Q - 1 = 0 \;\Rightarrow\; Q = \tfrac{3}{2}, \qquad \pi_A = \tfrac{9}{4}.
\]
Siembra cobra un precio $P = 1 - \epsilon$ y vende $Q = 3$, por lo que sus beneficios son $\pi_S = 3$.

\textbf{Caso 2: Agro impide la entrada con venta atada.} Para que la canasta de manzanas y naranjas sea preferida, la utilidad asociada (excedente: disposición máxima a pagar $V$ menos el precio) tiene que ser igual o mayor a la utilidad de comprar las manzanas por sí solas: $U_C \geq U_M$, es decir
\[
  V_N + V_M - P_C \geq V_M - P_M \quad\Longleftrightarrow\quad P_C \leq V_N + P_M .
\]
El precio de la canasta tiene que ser menor a la disposición a pagar por las naranjas más el precio de las manzanas. Debido a la competencia en precios en el mercado de las manzanas, Siembra cobrará a lo más su costo marginal, que es $0$. Por lo que el precio de la canasta tiene que ser menor a la disposición máxima a pagar por las naranjas (4). La demanda de la canasta es entonces $Q(P_C) = 16 - \frac{P_C^2}{2}$. Agro fija un precio de la canasta que cumple esta condición, de modo que el precio de las manzanas de Siembra debe ser $0$ para competir contra la canasta (la utilidad de comprar manzanas, para el consumidor, es la mayor posible dados los costos de producción de Siembra), y sus beneficios son $0$.

\emph{Nota:} la pauta original calculaba el precio óptimo de la canasta con la CPO de $\pi_A = (P_C - 1)\bigl(16 - \frac{P_C^2}{2}\bigr)$, obteniendo $P_C \approx 3{,}62$. Notar que el costo de producir la canasta es 2 (1 por cada bien); con ese costo la CPO lleva a $P_C = 4$, el borde del rango. En ambos casos se cumple la condición $P_C \leq 4$.

\textbf{En $T=1$ Siembra decide si entra.}
\begin{itemize}
  \item En el caso 1, donde Agro se acomoda: $\pi_S - F = 3 - 2 = 1 > 0$. En este caso sí habrá entrada.
  \item En el caso 2, donde Agro realiza una venta atada para disuadir la entrada: $\pi_S - F = 0 - 2 = -2 < 0$. No habrá entrada de Siembra en el mercado.
\end{itemize}
Por lo tanto, la venta atada sí disuade la entrada. Queda como propuesto ver si a Agro le conviene o no disuadir la entrada, comparando los beneficios asociados a cada caso.

\end{solution}
\end{parts}

%-------------------------------------------------------------------- Hotel Isla de Pascua
\needspace{8\baselineskip}
\tema{Bundling en un hotel (habitación y desayuno)}
\question Usted es el manager del hotel más grande de la Isla de Pascua y debe elegir los paquetes que ofrece a sus huéspedes. La tabla muestra la disposición a pagar de los turistas de tipo 1 y de tipo 2, tanto para una habitación (H) como para el desayuno (D). Además, se tiene información del número esperado de huéspedes de cada tipo y el costo marginal de proveer cada servicio.
\begin{center}
\begin{tabular}{lccc}
  Tipo de huésped & Habitación (H) & Desayuno (D) & N.º esperado de huéspedes \\ \hline
  Tipo 1 & 100 dólares & 5 dólares & 200 \\
  Tipo 2 & 60 dólares & 10 dólares & 800
\end{tabular}
\end{center}
Costo marginal: 40 dólares para H y 2 dólares para D.
\begin{parts}
\part Calcule el precio de la habitación $P_H$ y del servicio de desayuno $P_D$ que maximiza los beneficios del hotel. Encuentre los beneficios de vender cada servicio separadamente.
\needspace{5\baselineskip}
\begin{solution}
\textbf{Habitación.} El hotel puede poner el precio más bajo ($P_H = 60$) y vender a todos (1000 huéspedes): $\pi_H = (60 - 40)\cdot 1000 = 20\,000$. Si el hotel pone el precio más alto, $P_H = 100$, vende solo a los de tipo 1: $\pi_H = (100 - 40)\cdot 200 = 12\,000$. Por lo tanto, $P_H = 60$.

\textbf{Desayuno.} Si el precio es $P_D = 5$, vende a todos: $\pi_D = (5 - 2)\cdot 1000 = 3000$. Si el precio es $P_D = 10$, solo compran los de tipo 2: $\pi_D = (10 - 2)\cdot 800 = 6400$. El hotel obtiene mayor utilidad con $P_D = 10$.

Los beneficios totales son $\pi = 20\,000 + 6400 = 26\,400$ dólares.

\end{solution}

\part Suponga que el hotel ahora ofrece una habitación con desayuno incluido. Encuentre el precio del paquete $P_{HD}$ que maximiza los beneficios del hotel. Calcule los beneficios que obtiene el hotel al ofrecer el paquete.
\needspace{5\baselineskip}
\begin{solution}
Las disposiciones a pagar por el paquete son $105$ (tipo 1) y $70$ (tipo 2), y el costo marginal del paquete es $42$. El hotel puede vender el paquete a $P_{HD} = 70$ dólares, en cuyo caso le vende a todos y los beneficios son $\pi_{HD} = (70 - 42)\cdot 1000 = 28\,000$. Si el precio del paquete es $P_{HD} = 105$, vende solo a los que están dispuestos a pagar ese precio y los beneficios son $\pi_{HD} = (105 - 42)\cdot 200 = 12\,600$.

\end{solution}

\part ¿Qué le conviene más al hotel? ¿Ofrecer ventas atadas u ofrecer una habitación y el desayuno de manera separada?
\needspace{5\baselineskip}
\begin{solution}
La estrategia óptima es ofrecer la habitación con desayuno (paquete) a un precio $P_{HD} = 70$ dólares, con beneficios $\pi_{HD} = 28\,000 > 26\,400$.

\end{solution}
\end{parts}

%-------------------------------------------------------------------- Comente: bundling mixto vs puro
\needspace{8\baselineskip}
\tema{Comente: bundling puro vs bundling mixto}
\question \textbf{Comente.} Un monopolio multiproducto siempre preferirá realizar una venta atada por sobre una venta empaquetada (\emph{bundling} mixto) de sus productos, ya que, al limitar al máximo las opciones de las y los consumidores, es capaz de extraerles el máximo excedente posible.
\needspace{5\baselineskip}
\begin{solution}
Falso. Por lo general, el \emph{bundling} mixto tiene mayor capacidad de llevarse un mayor excedente de los consumidores: al vender en conjunto y por separado tiene más herramientas para extraer excedente. Considere también que al hacer \emph{bundling} mixto se relajan las restricciones de participación para ciertos individuos que no valoran mucho un bien relativo al otro. Por esto mismo, el \emph{bundling} puro puede llevar a perder consumidores; el uso de venta atada se justificaría, por ejemplo, para evitar la entrada de una firma.

\end{solution}

%-------------------------------------------------------------------- Comente: por qué ventas atadas
\needspace{8\baselineskip}
\tema{Comente: ¿por qué las empresas usan ventas atadas?}
\question \textbf{Comente.} Si ofrecer ventas atadas reduce el número de productos de una empresa, ¿por qué las empresas utilizan este mecanismo de precios?
\needspace{5\baselineskip}
\begin{solution}
Las empresas utilizan ventas atadas como un mecanismo de discriminación de precios, pues puede permitir aumentar utilidades al incorporar clientes que valoran mucho uno de los productos y por tanto compran el conjunto (y pagan por un bien que por separado no hubieran comprado). Cabe destacar que pueden perder clientes que no están dispuestos a comprar el paquete.

Por otra parte, las ventas atadas son un mecanismo de disuasión de entrada: la competencia que quiere entrar con un producto no es competitiva al no vender en paquete, y el incumbente fija precios de manera que a los clientes les convenga el paquete y el entrante se queda sin mercado.

\end{solution}

%-------------------------------------------------------------------- Venta atada Internet y TV
\needspace{8\baselineskip}
\tema{Venta atada, precios individuales y Doble Play (Internet y TV)}
\question Una empresa de telecomunicaciones vende dos productos: Internet y Televisión de Pago. Existen tres tipos de consumidores, cuyas máximas disposiciones a pagar por cada producto, para el consumidor representativo del grupo, son las que se indican a continuación. Existe igual número de consumidores en cada tipo. Considere un costo de producción igual a 0.
\begin{center}
\begin{tabular}{cccc}
  Consumidor & Internet & TV Pago & Suma \\ \hline
  A & 120 & 70 & 190 \\
  B & 90 & 100 & 190 \\
  C & 40 & 145 & 185
\end{tabular}
\end{center}
\begin{parts}
\part Si la empresa fija precio por productos individuales, ¿cuáles serían los precios fijados en cada uno de ellos?
\needspace{5\baselineskip}
\begin{solution}
Se venden los productos por separado, fijando precios que maximicen los beneficios.
\begin{itemize}
  \item \emph{Internet.} Hay tres opciones: que solo compren los de alta valoración ($P = 120$, beneficio $120\cdot 1$), que compren los de alta y media ($P = 90$, beneficio $90\cdot 2 = 180$), o que compren todos ($P = 40$, beneficio $40\cdot 3 = 120$). Se decide por $P_{Int} = 90$, sirviendo a $2/3$ del mercado.
  \item \emph{TV Pago.} $P = 145$: $145$; $P = 100$: $200$; $P = 70$: $210$. El precio debe apelar a todos los grupos, $P_{TV} = 70$.
\end{itemize}
Los beneficios son $\pi = \pi_{Int} + \pi_{TV} = 90\cdot 2 + 70\cdot 3 = 390$.

\end{solution}

\part La empresa decide realizar una venta atada, es decir, ofrecer los dos bienes en conjunto sin posibilidad de compra individual. ¿Cuál sería el precio que se cobraría por los dos bienes? ¿Le conviene esta estrategia a la empresa respecto de la venta individual?
\needspace{5\baselineskip}
\begin{solution}
Imponemos la restricción de participación para cada grupo: la valoración conjunta debe superar al precio para que haya venta. Las valoraciones conjuntas son $190$ (A), $190$ (B) y $185$ (C). Ante el \emph{trade-off} entre volumen y margen, toca decidir entre $190\cdot 2 = 380$ o $185\cdot 3 = 555$, donde claramente se impone $P = 185$ con beneficios $555$.

En comparación con el caso anterior hay mayores utilidades ($555 > 390$). Esto se debe a que los que valoraban mucho el Internet (grupo A) tienen que comprar también TV, mientras que los que valoraban más la TV (grupo C) tienen que comprar también Internet.

\end{solution}

\part La empresa decide vender los productos individuales, pero si un suscriptor desea comprar ambos productos (Doble Play), se le realiza un descuento respecto de la compra de ambos bienes por separado. Obtenga los tres precios óptimos que aplicaría la empresa (dos precios de productos separados y el precio Doble Play). Obtenga los beneficios totales de la empresa y compárelos con las dos preguntas anteriores. Comente.
\needspace{5\baselineskip}
\begin{solution}
Para planificar el \emph{bundling} mixto representamos todas las valoraciones por producto y paquete. Si alguien compra los bienes por separado, tienen que ser los que más lo valoran: $P_{Int} = 120$ (grupo A) y $P_{TV} = 145$ (grupo C). El precio del paquete tiene que considerar las restricciones de incentivo para que los de mayor valoración por TV e Internet sigan comprando los bienes separados.

El arte está en determinar el precio del paquete. Al fijar $P_{Int+TV} = 190 - \epsilon$, los grupos A y B compran el paquete: capturamos al consumidor intermedio B y sacamos mayor excedente de A, que no valoraba tanto la TV. Con este \emph{bundling} mixto tenemos $\pi = 190\cdot 2 + 145 = 525$.

Vemos que la venta atada pura ($555$) es mejor en este caso. Esto puede explicarse por la gran diferencia en valoración entre bienes y grupos, por lo que el paquete tiende a sacar mayor excedente al juntar el bien muy valorado con el menos valorado. Adicionalmente, la valoración conjunta en los tres grupos es similar, por lo que bajar 5 unidades el precio para capturar al restante $1/3$ es altamente beneficioso. (Notar que el \emph{bundling} mixto óptimo nunca puede ser peor que el puro, pues puede replicarlo fijando precios individuales muy altos; aquí el óptimo mixto coincide con no vender por separado.)

\end{solution}
\end{parts}

%-------------------------------------------------------------------- Comente: menú restaurante
\needspace{8\baselineskip}
\tema{Comente: menú y platos a la carta}
\question \textbf{Comente.} Cuando una firma como un restaurante ofrece un menú que incluye varios platos por un precio único por el conjunto, pero también ofrece los platos por separado a la carta, entonces todos los consumidores siempre estarán mejor que si no se ofrece el menú, ya que siempre es mejor tener más opciones que menos.
\needspace{5\baselineskip}
\begin{solution}
Falso. El menú de opciones le permite al restaurante subir los precios de los platos a la carta, extrayendo más excedente de los consumidores que valoran mucho alguno de estos platos en forma individual pero no los otros del menú, sin perder la demanda de los consumidores con gustos intermedios por todos los platos. Así, hay algunos consumidores que ganan al ofrecerse el menú (los de gustos intermedios) y otros que pierden (los de gustos extremos por algún plato en particular).

\end{solution}

%-------------------------------------------------------------------- Triple pack
\needspace{8\baselineskip}
\tema{Triple pack vs servicios por separado (Internet, TV y telefonía)}
\question Una empresa de telecomunicaciones ofrece tres servicios: Internet, Televisión y Telefonía. Existen tres tipos de consumidoras, clasificadas como A, B y C. La máxima disposición a pagar por categoría de servicio y según tipo de consumidora es la siguiente:
\begin{center}
\begin{tabular}{cccc}
  Consumidoras & Internet & Televisión & Telefonía \\ \hline
  A & 13 & 6 & 3 \\
  B & 7 & 12 & 16 \\
  C & 6 & 25 & 11
\end{tabular}
\end{center}
Suponga que hay un mismo número de consumidoras por cada categoría. Todos los costos fijos y variables son cero.
\begin{parts}
\part Si la empresa vende solamente un único plan que incluye todos los servicios, ¿cuál sería el precio que cobra por el triple pack? ¿Cuánto sería el beneficio obtenido?
\needspace{5\baselineskip}
\begin{solution}
Las valoraciones del triple pack son $22$ (A), $35$ (B) y $42$ (C). Normalizando el número de consumidoras de cada tipo a 1:
\[
  P = 42 \Rightarrow 1\text{ compradora}: 42; \qquad
  P = 35 \Rightarrow 2: 70; \qquad
  P = 22 \Rightarrow 3: 66 .
\]
Por lo tanto, la firma fijará $P = 35$ para el triple pack y obtendrá beneficios de $70$.

\end{solution}

\part La empresa está considerando vender los servicios por separado. ¿Cuál sería el precio cobrado por cada servicio? ¿Le conviene cambiar su estrategia de precios respecto al plan único? Explique.
\needspace{5\baselineskip}
\begin{solution}
Analizamos precios, cantidades y beneficios de cada servicio:
\begin{itemize}
  \item \emph{Internet:} $P=13 \to 13$; $P=7 \to 14$; $P=6 \to 18$. Conviene $P = 6$ (beneficio $18$).
  \item \emph{Televisión:} $P=25 \to 25$; $P=12 \to 24$; $P=6 \to 18$. Conviene $P = 25$ (beneficio $25$).
  \item \emph{Telefonía:} $P=16 \to 16$; $P=11 \to 22$; $P=3 \to 9$. Conviene $P = 11$ (beneficio $22$).
\end{itemize}
Sumando los beneficios, el total es $18 + 25 + 22 = 65$. Al compararlo con el beneficio del triple pack ($70$), concluimos que no le conviene a la empresa cambiar su estrategia de precios. La razón es que al ofrecer un paquete único, la firma puede capturar mayor disposición a pagar de consumidores con preferencias diversas, mientras que al desagregar los servicios se pierde esta ventaja.

\end{solution}
\end{parts}

%-------------------------------------------------------------------- HP impresoras y toner
\needspace{8\baselineskip}
\tema{Bundling puro y mixto: impresoras y cartuchos de tóner (HP)}
\question HP vende dos productos: impresoras y cartuchos de tóner. Hay tres consumidores en el mercado. El consumidor 1 necesita imprimir con mucha frecuencia, por lo que valora mucho ambos bienes; el consumidor 2 necesita imprimir con menos frecuencia, por lo que valora menos ambos bienes; y el consumidor 3 ya tiene impresora, por lo que solo valora el tóner, y lo valora mucho porque ya ha realizado el gasto en la impresora. Las valoraciones de los consumidores aparecen en la siguiente tabla:
\begin{center}
\begin{tabular}{lcc}
  & Impresora & Tóner \\ \hline
  Consumidor tipo 1 & 8 & 4 \\
  Consumidor tipo 2 & 6 & 3 \\
  Consumidor tipo 3 & 0 & 8
\end{tabular}
\end{center}
Para simplificar, el costo marginal de producir ambos bienes es cero. Calcule los precios óptimos y los beneficios de HP en las siguientes situaciones:
\begin{parts}
\part HP vende ambos bienes de forma separada.
\needspace{5\baselineskip}
\begin{solution}
Hay dos posibilidades para el precio de la impresora: $P(I) = 8$ (solo compra el consumidor 1; $\pi = 8$) o $P(I) = 6$ (compran los consumidores 1 y 2; $\pi = 12$). Hay tres posibilidades para el precio del tóner: $P(T) = 8$ (solo compra el consumidor 3; $\pi = 8$), $P(T) = 4$ (compran 1 y 3; $\pi = 8$) o $P(T) = 3$ (compran los tres; $\pi = 9$). Luego, los precios óptimos son $P(I) = 6$ y $P(T) = 3$, y las utilidades totales son $\pi = 12 + 9 = 21$.

\end{solution}

\part HP aplica ventas atadas y solo vende los dos bienes de forma conjunta. ¿Qué preferiría HP: vender ambos bienes separadamente o de forma conjunta?
\needspace{5\baselineskip}
\begin{solution}
Las valoraciones del paquete son $12$, $9$ y $8$ (consumidores 1, 2 y 3). Hay tres posibilidades: $P(I+T) = 12$ (compra solo el 1; $\pi = 12$), $P(I+T) = 9$ (compran 1 y 2; $\pi = 18$), $P(I+T) = 8$ (compran los tres; $\pi = 24$). Por lo tanto, el precio óptimo es $P(I+T) = 8$ y las utilidades son $\pi = 24$. HP preferirá vender los bienes en conjunto, ya que las utilidades son mayores ($24 > 21$).

\end{solution}

\part HP también puede aplicar empaquetamiento y vender los bienes tanto de forma conjunta como separada. ¿Cuál será la elección final de HP? ¿Vender los bienes de forma separada, aplicar ventas atadas o empaquetamiento?
\needspace{5\baselineskip}
\begin{solution}
Con $P(I+T) = 9$, $P(I) = 8$ y $P(T) = 8$: los consumidores 1 y 2 compran el paquete (el consumidor 1 obtiene un excedente de $3$ con el paquete y $0$ con la impresora sola; el consumidor 2 no puede comprar la impresora a $8$) y el consumidor 3 compra el tóner. Las utilidades son $\pi = 9 + 9 + 8 = 26$. (Cualquier precio de la impresora entre $6$ y $8$ funciona; debe ser al menos $6$, porque si fuera menor el consumidor 2 compraría solo la impresora, que le da un excedente positivo en lugar de cero.)

Notar que el esquema $P(I+T) = 12$, $P(I) = 6$, $P(T) = 8$ no es compatible en incentivos: con $P(I) = 6$, el consumidor 1 prefiere comprar solo la impresora (excedente $2$) al paquete (excedente $0$), por lo que HP obtendría solo $6 + 6 + 8 = 20$.

Por consiguiente, la elección final de HP será el empaquetamiento (vender por separado y en conjunto), ya que las utilidades ($\pi = 26$) son mayores que si vendiera solo por separado ($\pi = 21$) o ventas atadas ($\pi = 24$).

\end{solution}
\end{parts}
